\section{本章小结}
\label{sec:graph-incidence-summary}

给定参考方向后，图的带符号关联矩阵$\bm{B}\in\mathbb{R}^{n\times m}$在第$e$列（对应有向边$e\colon u\to v$）于尾端$u$记$+1$、头端$v$记$-1$，把顶点函数空间$\mathbb{R}^V$与边空间$\mathbb{R}^E$连接成一对伴随映射：转置矩阵$\bm{B}^\top\colon\mathbb{R}^V\to\mathbb{R}^E$将顶点函数$\bm{f}$映射为沿边差分向量$\bm{d}=\bm{B}^\top\bm{f}$，矩阵本身$\bm{B}\colon\mathbb{R}^E\to\mathbb{R}^V$将边流量向量$\bm{q}$映射为各顶点的净流出向量$\bm{B}\bm{q}$，二者满足伴随恒等式$\langle\bm{B}^\top\bm{f},\bm{q}\rangle_E=\langle\bm{f},\bm{B}\bm{q}\rangle_V$。

在边空间$\mathbb{R}^E$中，每个顶点处净流出量均为零的边流量构成环空间$\mathcal{C}(\mathcal{G})=\ker\bm{B}$，能够表示为顶点函数沿边差分向量的边向量构成割空间$\mathcal{D}(\mathcal{G})=\operatorname{im}\bm{B}^\top$。由伴随恒等式，环空间与割空间正交互补，将边空间正交直和分解为$\mathbb{R}^E=\ker\bm{B}\oplus\operatorname{im}\bm{B}^\top$。

若图有$n$个顶点、$m$条边和$c$个连通分量，则$\ker\bm{B}^\top$由$c$个连通分量的指示向量张成且线性无关，维数为$c$；由秩零度定理，关联矩阵的秩与割空间维数均为$n-c$，环空间维数为$m-n+c$。

选定一棵生成森林$F$后，向$F$中逐一添入$m-n+c$条非树边得到环空间的基本环基（在非树边坐标上构成单位阵），从$F$中逐一删去$n-c$条树边得到割空间的基本割基（在树边坐标上构成单位阵）。

在组合优化中，若整数矩阵的任意阶方子式均属于$\{0,1,-1\}$，则称之为全幺模矩阵。在标准型线性规划中，全幺模系数矩阵与整数右端项使可行域的每个可逆基矩阵行列式均为$\pm1$，由克莱姆法则保证所有极点（基本可行解）均为整数向量，当问题可行且最优值有限时，便存在整数最优解。

按约束矩阵类型划分：以带符号关联矩阵$\bm{B}$为约束的有向网络流家族（最小费用流、最大流、最短路），因$\bm{B}$每列恰有一正一负、对任意无自环图恒全幺模，故在整数供需与容量、问题可行且最优值有限时存在整数最优流；最短路的单位流模型还须排除负费用环。

以无符号关联矩阵$\bm{N}$为约束的配对与覆盖家族（最大匹配、指派问题、最小顶点覆盖），因$\bm{N}$每列有两个$+1$、当且仅当图为不含奇环的二部图时全幺模（奇环子矩阵的排列展开给出行列式$1+1=2$），故在二部图上保证存在$\{0,1\}$整数最优解，而在含奇环图上会出现分数极点。

后续章节将频繁使用本章建立的关联矩阵框架。将沿边差分算子$\bm{B}^\top$、非负边权对角矩阵$\bm{C}$与净流出汇总算子$\bm{B}$依次复合，得到顶点空间上的对称半正定算子$\bm{L}=\bm{B}\bm{C}\bm{B}^\top$——这正是第~\ref{chap:graph-laplacian-matrix}章的组合拉普拉斯矩阵，并在第~\ref{chap:graph-electrical-networks}章的电网络方程中把电势差、支路电导与基尔霍夫电流守恒融为一体。
